\chapter{Application Workflows}
\label{ch:application_workflow}

This chapter details the operational workflows of RGUKT Connect, outlining the steps, endpoints, and data transformations for key processes.

\section{Registration \& OTP Verification Workflow}
Establishing a new user account requires dual verification: university email validation and OTP challenge.
\begin{enumerate}
    \item \textbf{Identity Submission:} The user enters their name, ID number, and email.
    \item \textbf{Backend Validation:} The frontend sends a request to:
    \begin{lstlisting}[language=bash]
    POST /api/auth/register/send-otp?email={email}&idNumber={idNumber}
    \end{lstlisting}
    The backend verifies the email domain is `@rgukt.ac.in`, the ID is 7 characters, and that neither the email nor the ID is already registered.
    \item \textbf{OTP Generation:} The backend generates a 6-digit OTP, saves it in an in-memory map (`registrationOtpStorage`) with a 5-minute expiry, and sends it via SMTP.
    \item \textbf{Bypass Check:} If the email starts with `test`, the system skips SMTP and sets the OTP to `123456`.
    \item \textbf{Submission:} The user submits their registration details, password, and the OTP to:
    \begin{lstlisting}[language=bash]
    POST /api/auth/register
    \end{lstlisting}
    \item \textbf{Commit Phase:} The backend validates the OTP. If valid, it hashes the password with BCrypt, saves a new `User` record, creates a default `UserDetails` entry (setting default branch to `CSE` and phone to `9876543210`), and deletes the OTP from the in-memory map.
\end{enumerate}

\section{Authentication \& JWT Lifecycle}
Logging in establishes a stateless session:
\begin{enumerate}
    \item \textbf{Credentials Submission:} The user submits their email and password to:
    \begin{lstlisting}[language=bash]
    POST /api/auth/login
    \end{lstlisting}
    \item \textbf{Credential Verification:} The backend verifies the email domain. It retrieves the user, validates the BCrypt password hash, and returns an error if authentication fails.
    \item \textbf{Token Issuance:} The system generates a JWT access token containing the user's email, ID number, and role, signed with the backend secret.
    \item \textbf{Client Cache:} The backend returns the JWT token and basic profile metadata. The React client saves this data as a JSON string in LocalStorage under the key `userSession`.
    \item \textbf{Authorized Requests:} The client includes the token in the `Authorization` header (`Bearer <token>`) for all subsequent API requests.
\end{enumerate}

\section{Job Posting \& Referral Workflow}
Posting a job opportunity requires specific role privileges:
\begin{enumerate}
    \item \textbf{Form Submission:} The user fills out the job details form.
    \item \textbf{API Request:} The client sends the payload with the authorization token to:
    \begin{lstlisting}[language=bash]
    POST /api/jobs
    \end{lstlisting}
    \item \textbf{Role Validation:} The backend extracts the user's role from the JWT. If the role is not `ALUMNI` or `ADMIN`, the backend rejects the request with an HTTP 403 status code.
    \item \textbf{DB Write:} If authorized, the backend creates a `Job` record linked to the posting user.
\end{enumerate}

\section{Connection \& Real-Time Notification Workflow}
Establishing a connection updates the network graph and triggers notifications:
\begin{enumerate}
    \item \textbf{Connection Request:} User A clicks "Connect" on User B's profile. The client sends a request to:
    \begin{lstlisting}[language=bash]
    POST /api/connections/request/{userId}
    \end{lstlisting}
    \item \textbf{Database Entry:} The backend creates a `Connection` record with a status of `PENDING`.
    \item \textbf{Notification Trigger:} The backend creates a `Notification` entry with the type \texttt{CONNECTION\_REQUEST} and the ID of the connection request.
    \item \textbf{Real-Time Push:} The backend pushes the notification over WebSockets to the recipient:
    \begin{lstlisting}[language=bash]
    /user/{recipientId}/queue/notifications
    \end{lstlisting}
    \item \textbf{Acceptance:} User B clicks "Accept". The client sends a request to:
    \begin{lstlisting}[language=bash]
    PUT /api/connections/accept/{requestId}
    \end{lstlisting}
    The backend updates the connection status to `ACCEPTED` and pushes a \texttt{CONNECTION\_ACCEPTED} notification back to User A.
\end{enumerate}

\section{Real-Time Messaging Workflow}
Private messaging uses WebSockets to support real-time delivery:
\begin{enumerate}
    \item \textbf{STOMP Handshake:} The client establishes a WebSocket connection:
    \begin{lstlisting}[language=bash]
    ws://<host>/ws-chat/websocket
    \end{lstlisting}
    A custom channel interceptor validates the JWT token in the `Authorization` header.
    \item \textbf{Subscription:} The client subscribes to:
    \begin{lstlisting}[language=bash]
    /user/queue/messages
    \end{lstlisting}
    \item \textbf{Message Dispatch:} When sending a message, the client posts the JSON payload to `/api/chat/send` via an HTTP POST request.
    \item \textbf{Persistence:} The backend saves the message to the database and checks if User B is currently online.
    \item \textbf{WebSocket Delivery:} The system pushes the message to User B:
    \begin{lstlisting}[language=bash]
    /user/{receiverId}/queue/messages
    \end{lstlisting}
\end{enumerate}
